\documentclass[journal]{IEEEtran}

\usepackage{amsmath,amssymb}
\usepackage{graphicx}
\usepackage{xcolor}
\usepackage{array}
\usepackage{cite}

\newcommand{\needref}[1]{\textcolor{red}{[#1]}}
\newcommand{\RR}{\mathrm{RR}}
\newcommand{\RI}{\mathrm{RI}}
\newcommand{\IR}{\mathrm{IR}}
\newcommand{\II}{\mathrm{II}}
\newcommand{\diag}{\mathrm{diag}}
\newcommand{\sq}{\mathrm{sq}}
\graphicspath{{../logs/}}

\title{Lifted Rectangular AC Optimal Power Flow via Linearized Augmented Lagrangian and Quantum Hamiltonian Descent}

\author{Author Name(s)
\thanks{The authors are with the Department/Institute Name, University/Organization Name. E-mail: author@example.com.}
\thanks{This manuscript is a working draft. Red bracketed text indicates locations where references should be inserted before submission.}}

\begin{document}

\maketitle

\begin{abstract}
AC optimal power flow (ACOPF) is a nonconvex network optimization problem whose equality constraints are difficult to encode directly in Ising or QUBO forms used by quantum and quantum-inspired optimizers \needref{ACOPF nonconvexity and OPF survey references}. This paper develops a lifted rectangular-coordinate ACOPF framework in which active and reactive power balance, branch-flow definitions, voltage-magnitude equations, thermal-flow auxiliaries, and reference-bus conditions are collected in an augmented Lagrangian. To preserve a quadratic subproblem structure, nonlinear lifted-consistency and thermal auxiliary constraints are locally linearized at each outer iteration, while affine constraints in the lifted variables are retained exactly. The resulting box-constrained quadratic subproblems are solved with a Quantum Hamiltonian Descent (QHD) implementation using a Simulated Bifurcation backend, unary discretization, spin decoding, and local continuous refinement. Preliminary Simbi-only comparisons on 5-, 9-, and 14-bus systems are reported against standard-answer files in the repository. The 14-bus run reaches a maximum residual of $1.05\times 10^{-4}$, while the selected 5- and 9-bus runs expose the current feasibility gap that remains after QHD/Simbi sampling and local refinement. The current results are intended as a modeling and algorithmic baseline rather than a claim of quantum advantage \needref{QHD, quantum annealing, and simulated bifurcation references}.
\end{abstract}

\begin{IEEEkeywords}
AC optimal power flow, augmented Lagrangian method, rectangular coordinates, lifted formulation, QUBO, Ising optimization, Quantum Hamiltonian Descent, simulated bifurcation.
\end{IEEEkeywords}

\section{Introduction}
\IEEEPARstart{A}{C} optimal power flow (ACOPF) is a central optimization problem in electric power system operation. It minimizes generation cost subject to Kirchhoff-law power balance, nonlinear AC branch-flow equations, generator limits, voltage limits, and thermal line limits. The nonconvexity of ACOPF makes large-scale and globally reliable solution challenging, motivating convex relaxations, nonlinear programming algorithms, decomposition methods, and learning-assisted or physics-informed variants \needref{classical ACOPF algorithms and relaxations}.

Quantum and quantum-inspired optimization methods offer another possible computational route: if an OPF subproblem can be written as a quadratic Hamiltonian or QUBO, it can be submitted to annealing-style samplers, simulated bifurcation engines, or Hamiltonian-descent solvers \needref{QUBO/Ising optimization and power-system quantum optimization references}. The difficulty is that the native ACOPF equations are not immediately quadratic in a convenient binary encoding. Direct penalty formulations can become high order, poorly scaled, or too constraint-heavy for current hardware and simulators. A practical formulation must therefore separate the power-system modeling issue from the solver-encoding issue.

This work studies that separation. We use a lifted rectangular-coordinate ACOPF model in which voltage products are represented by auxiliary variables. In this lifted space, many network equations become affine. The remaining nonlinear consistency equations are handled through a linearized augmented Lagrangian method (ALM), so that each outer iteration produces a quadratic, box-constrained subproblem. This subproblem is then passed to a QHD-style interface and solved with the local Simulated Bifurcation backend implemented in the adjusted \texttt{qhdopt} code.

The main contributions are as follows.
\begin{itemize}
    \item A sparse lifted rectangular ACOPF formulation that keeps only the voltage-product variables required by branch-flow equations and nonzero admittance entries.
    \item A linearized ALM construction that preserves a quadratic subproblem by linearizing only nonlinear residual blocks, while retaining affine physical equations exactly.
    \item A QHD/Simbi implementation path that maps each boxed quadratic subproblem to a unary-encoded spin model and decodes samples back to continuous OPF variables.
    \item Preliminary numerical evidence, based only on Simbi runs, on 5-, 9-, and 14-bus systems, including comparison with repository standard-answer files.
\end{itemize}

\section{Related Work}
Classical ACOPF methods include nonlinear interior-point and sequential quadratic programming solvers, semidefinite and second-order-cone relaxations, moment relaxations, convex envelopes, and decomposition methods \needref{interior-point OPF, SDP/SOC OPF relaxations, moment relaxation references}. These methods are mature and remain the appropriate baseline for operational studies. The present paper does not replace those solvers; instead, it constructs an OPF subproblem family compatible with Hamiltonian or QUBO optimizers.

The second relevant line of work concerns quantum annealing, gate-model variational optimization, Hamiltonian-descent formulations, and simulated bifurcation methods for quadratic unconstrained binary optimization \needref{quantum annealing, QAOA, QHD, simulated bifurcation references}. Power-system applications of these methods usually require either binary expansion of continuous variables, discretized OPF approximations, or penalty methods for equality constraints \needref{quantum optimization for OPF and unit commitment references}. This paper differs by maintaining a continuous lifted AC model at the outer level, then producing a quadratic encoded subproblem at each ALM iteration.

\section{Lifted Rectangular ACOPF Model}
\subsection{Network Sets and Data}
Let $\mathcal N$ denote the set of buses, $\mathcal E$ the set of undirected transmission lines, and $\mathcal A$ the set of directed arcs obtained by including both directions of each line. Let $\mathcal G_i$ denote the set of generators at bus $i$. The bus admittance matrix is $Y=G+jB$, and the series admittance of directed arc $a=(i,j)$ is $g_a+j b_a$. Active and reactive loads are $P_i^D$ and $Q_i^D$. Each generator $g$ has limits $P_g^{\min}$, $P_g^{\max}$, $Q_g^{\min}$, $Q_g^{\max}$ and quadratic cost
\begin{equation}
    C_g(P_g)=a_g P_g^2+b_g P_g+c_g .
\end{equation}

\subsection{Decision Variables}
The primal vector contains generator powers, lifted voltage products, voltage-magnitude auxiliaries, branch-flow auxiliaries, and squared apparent-power auxiliaries:
\begin{equation}
    x=\left(P_G,Q_G,W_{\diag},W_{\mathrm{cross}},V^{\sq},P_{\mathcal A},Q_{\mathcal A},S_{\mathcal A}^{\sq}\right).
\end{equation}
For each bus $i$, the diagonal lifted terms are
\begin{equation}
    W_{ii}^{\RR}=(V_i^R)^2,\quad
    W_{ii}^{\RI}=V_i^R V_i^I,\quad
    W_{ii}^{\II}=(V_i^I)^2 .
\end{equation}
For each sparse lifted pair $(i,j)$ required by the physical network, four cross terms are used:
\begin{equation}
\begin{aligned}
    W_{ij}^{\RR}=V_i^R V_j^R,\quad
    W_{ij}^{\RI}=V_i^R V_j^I,\\
    W_{ij}^{\IR}=V_i^I V_j^R,\quad
    W_{ij}^{\II}=V_i^I V_j^I .
\end{aligned}
\end{equation}
In the implementation, this pair set is built from branch-flow pairs and nonzero off-diagonal admittance entries. Hence the lifted model is sparse rather than all-to-all.

The remaining feasible set is a box
\begin{equation}
    \mathcal X=\{x:\underline{x}\leq x\leq \overline{x}\},
\end{equation}
where the bounds come from generator limits, voltage-magnitude limits, branch MVA ratings, and conservative bounds on lifted voltage products and branch-flow variables.

\subsection{Objective}
The generation-cost objective is
\begin{equation}
    f(x)=\sum_{g\in\mathcal G} C_g(P_g).
\end{equation}
The lifted variables enter the objective only through the augmented terms described below.

\subsection{Equality Residuals}
All physical equations are stacked into a residual vector $h(x)$. For bus $i$, the active-power balance residual is
\begin{align}
c_i^P(x)={}&\sum_{g\in\mathcal G_i}P_g-P_i^D
    -G_{ii}(W_{ii}^{\RR}+W_{ii}^{\II}) \nonumber\\
&-\sum_{j\in\mathcal N_i}
\bigl(G_{ij}W_{ij}^{\RR}-B_{ij}W_{ij}^{\RI} \nonumber\\
&\hspace{3.0em}
      +G_{ij}W_{ij}^{\II}+B_{ij}W_{ij}^{\IR}\bigr),
\end{align}
where $\mathcal N_i$ contains neighbors with nonzero admittance terms. The reactive-power balance residual is
\begin{align}
c_i^Q(x)={}&\sum_{g\in\mathcal G_i}Q_g-Q_i^D
    +B_{ii}(W_{ii}^{\RR}+W_{ii}^{\II}) \nonumber\\
&-\sum_{j\in\mathcal N_i}
\bigl(G_{ij}W_{ij}^{\IR}-B_{ij}W_{ij}^{\II} \nonumber\\
&\hspace{3.0em}
      -G_{ij}W_{ij}^{\RI}-B_{ij}W_{ij}^{\RR}\bigr).
\end{align}

For directed arc $a=(i,j)$, branch active and reactive flow definitions are
\begin{align}
c_a^{P,f}(x)={}&P_a-
g_a\bigl((W_{ii}^{\RR}+W_{ii}^{\II}) \nonumber\\
&-(W_{ij}^{\RR}+W_{ij}^{\II})\bigr)
-b_a(W_{ij}^{\RI}-W_{ij}^{\IR}),\\
c_a^{Q,f}(x)={}&Q_a-
b_a\bigl((W_{ij}^{\RR}+W_{ij}^{\II}) \nonumber\\
&-(W_{ii}^{\RR}+W_{ii}^{\II})\bigr)
-g_a(W_{ij}^{\RI}-W_{ij}^{\IR}).
\end{align}
The voltage-magnitude and branch thermal auxiliary residuals are
\begin{align}
c_i^V(x)&=V_i^{\sq}-(W_{ii}^{\RR}+W_{ii}^{\II}),\\
c_a^S(x)&=S_a^{\sq}-(P_a^2+Q_a^2).
\end{align}
The lifted consistency residuals are
\begin{align}
c_i^{W,d}(x)&=W_{ii}^{\RR}W_{ii}^{\II}-(W_{ii}^{\RI})^2,\\
c_{ij}^{W,c}(x)&=(W_{ij}^{\RR}+W_{ij}^{\II})^2
 +(W_{ij}^{\RI}-W_{ij}^{\IR})^2
 -V_i^{\sq}V_j^{\sq}.
\end{align}
Finally, with reference bus $r$, the current implementation enforces
\begin{equation}
    W_{rr}^{\II}=0,\quad W_{rr}^{\RI}=0,\quad
    W_{rj}^{\II}=0,\quad W_{rj}^{\IR}=0
\end{equation}
for incident lifted pairs $(r,j)$. The full residual vector is ordered as
\begin{equation}
h(x)=\left[
c^P,c^Q,c^{P,f},c^{Q,f},c^V,c^S,c^{W,d},c^{W,c},c^{\mathrm{ref}}
\right]^T,
\end{equation}
where each entry denotes the corresponding stacked residual block.

\section{Linearized Augmented Lagrangian}
The fully augmented problem would minimize
\begin{equation}
    \mathcal L_\rho(x,\lambda)=
    f(x)+\lambda^T h(x)+\frac{\rho}{2}\|h(x)\|_2^2
\end{equation}
over $x\in\mathcal X$. However, directly squaring nonlinear residuals would generate high-order polynomial terms that are not compatible with a quadratic Hamiltonian. The implementation therefore separates residuals into affine blocks $h_L(x)$ and nonlinear blocks $h_N(x)$. In the present model, the nonlinear blocks are $c^S$, $c^{W,d}$, and $c^{W,c}$.

At outer iteration $k$, nonlinear blocks are replaced by their first-order approximation
\begin{equation}
    \widetilde h_N^{\,k}(x)=h_N(x^k)+J_N(x^k)(x-x^k),
\end{equation}
while affine residuals are kept exact. Let
\begin{equation}
    \widetilde h_k(x)=\left[h_L(x),\widetilde h_N^{\,k}(x)\right]^T .
\end{equation}
The QHD-compatible subproblem is
\begin{equation}
\label{eq:linearized-alm-subproblem}
\begin{aligned}
x^{k+1}\in\arg\min_{x\in\mathcal X}\;&
 f(x)+(\lambda^k)^T\widetilde h_k(x)\\
&+\frac{\rho_k}{2}\|\widetilde h_k(x)\|_2^2
 +\frac{\mu}{2}\|x-x^k\|_2^2 .
\end{aligned}
\end{equation}
Because $f(x)$ is quadratic and $\widetilde h_k(x)$ is affine, \eqref{eq:linearized-alm-subproblem} is a box-constrained quadratic program. After solving the subproblem, the dual variables are updated using the true nonlinear residual:
\begin{equation}
    \lambda^{k+1}=\lambda^k+\alpha_k h(x^{k+1}).
\end{equation}
The implementation includes adaptive adjustment of $\alpha_k$ and $\rho_k$: stable residual decrease can slightly increase $\alpha_k$, residual worsening decreases $\alpha_k$, repeated worsening increases $\rho_k$, and plateau behavior increases $\rho_k$ while damping $\alpha_k$.

\section{QHD and Simulated Bifurcation Subproblem Solver}
Each quadratic subproblem is submitted through the \texttt{QHD.SymPy} interface. Bounds are handled by an affine map
\begin{equation}
    x=\underline{x}+D z,\qquad z\in[0,1]^n,
\end{equation}
where $D=\mathrm{diag}(\overline{x}-\underline{x})$. With unary resolution $r$, each scalar $z_i$ is represented by $r$ binary variables. The decoded value is the unary average when the bitstring is monotone; otherwise, the implementation projects to a soft unary representation before decoding.

The expanded quadratic polynomial is decomposed into univariate and bivariate terms. These terms are compiled to an Ising energy
\begin{equation}
    E(s)=\sum_{i<j}J_{ij}s_i s_j+\sum_i h_i s_i,\qquad s_i\in\{-1,+1\}.
\end{equation}
The Simulated Bifurcation backend solves the equivalent dense quadratic spin polynomial
\begin{equation}
    \min_s\; s^T M s+v^T s,
\end{equation}
with $M=\frac{1}{2}J$ and $v=h$. The sampled spins are converted to bitstrings, decoded to continuous subproblem variables, and optionally refined by a continuous local optimizer. In the reported runs, the Simbi backend was used with unary embedding, resolution $r=14$, $128$ shots, $4096$ agents, $1300$ maximum SB steps, and local TNC post-processing.

\section{Algorithmic Summary}
The implemented procedure is:
\begin{enumerate}
    \item Build the sparse lifted rectangular ACOPF model and initialize $x^0$ from load sharing, voltage magnitudes, and zero branch-flow auxiliaries.
    \item Initialize $\lambda^0=0$, $\rho_0$, $\alpha_0$, and the proximal weight $\mu$.
    \item At each outer iteration, construct $\widetilde h_k(x)$ by linearizing only nonlinear residual blocks at $x^k$.
    \item Form the quadratic ALM subproblem in \eqref{eq:linearized-alm-subproblem}.
    \item Solve the boxed quadratic subproblem with QHD/Simbi and local refinement.
    \item Evaluate the true residual $h(x^{k+1})$, update $\lambda^{k+1}$, and adapt $\rho_k,\alpha_k$.
    \item Stop when the maximum absolute residual is below tolerance or when the early-stop patience criterion is reached; retain the best residual iterate.
\end{enumerate}

\section{Numerical Study}
\subsection{Test Cases and Encoded Sizes}
The current numerical study uses only Simbi backend runs. The standard answers are taken from \texttt{5bus-answer.txt}, \texttt{9bus-answer.txt}, and \texttt{14bus-answer.txt}. Table~\ref{tab:size} summarizes the lifted model sizes. The logical spin count is the number of continuous variables multiplied by the unary resolution $r=14$; it is not a hardware physical-qubit count.

\begin{table}[!t]
\caption{Lifted ACOPF Size and Unary-Encoded Spin Dimension}
\label{tab:size}
\centering
\begin{tabular}{c|c|c|c|c|c|c}
\hline
Case & Buses & Lines & Gen. & Variables & Residuals & Spins \\
\hline
5-bus  & 5  & 6  & 3 & 86  & 70  & 1204 \\
9-bus  & 9  & 9  & 3 & 132 & 103 & 1848 \\
14-bus & 14 & 20 & 5 & 266 & 202 & 3724 \\
\hline
\end{tabular}
\end{table}

\subsection{Residual and Objective Summary}
Table~\ref{tab:simbi-results} reports the best iterate by $\ell_2$ residual from the selected logs. The objective gap is computed as $100(f_{\mathrm{QHD}}-f_{\mathrm{std}})/f_{\mathrm{std}}$. A negative objective gap should not be interpreted as a better OPF solution when residuals are nonzero; it indicates that the sampled point is still not exactly feasible. The ``feasible'' flag in the raw log remains false because the checker uses a strict equality tolerance, so the table reports residual magnitudes directly.

\begin{table*}[!t]
\caption{Preliminary Simbi-Only Results Compared with Standard Answer Files}
\label{tab:simbi-results}
\centering
\begin{tabular}{c|c|c|c|c|c|c|c}
\hline
Case & Standard file & $f_{\mathrm{std}}$ & Selected QHD log & Best iter. & $f_{\mathrm{QHD}}$ & Gap & $\max |h(x)|$ \\
\hline
5-bus  & \texttt{5bus-answer.txt}  & 9.195047 & \texttt{Buses-5\_06-07-2026\_00-05-09}  & 189 & 9.162802 & -0.351\% & $2.899{\times}10^{-3}$ \\
9-bus  & \texttt{9bus-answer.txt}  & 4.098132 & \texttt{Buses-9\_06-08-2026\_19-48-09}  & 17  & 4.045797 & -1.277\% & $7.509{\times}10^{-3}$ \\
14-bus & \texttt{14bus-answer.txt} & 4.875262 & \texttt{Buses-14\_06-08-2026\_11-02-24} & 86  & 4.910898 & 0.731\% & $1.053{\times}10^{-4}$ \\
\hline
\end{tabular}
\end{table*}

\begin{table*}[!t]
\caption{Euclidean Distance to Standard Answers for Best-Residual Iterates}
\label{tab:load-results}
\centering
\begin{tabular}{c|c|c|c|c|c|c}
\hline
Case & Iter. & Bus L2 & Branch L2 & Combined L2 & Objective diff. & Load supplied \\
\hline
5-bus  & 189 & 0.118648 & 0.114625 & 0.164974 & -0.032244 & 99.891\% \\
9-bus  & 17  & 0.647244 & 1.129135 & 1.301488 & -0.052335 & 99.316\% \\
14-bus & 86  & 0.670897 & 1.011694 & 1.213930 & 0.035636  & 99.984\% \\
\hline
\end{tabular}
\end{table*}

\begin{table}[!t]
\caption{Closest-to-Standard Iterates by Combined L2 Distance}
\label{tab:closest-standard}
\centering
\begin{tabular}{c|c|c|c}
\hline
Case & Iter. & Combined L2 & Objective diff. \\
\hline
5-bus  & 134 & 0.026343 & -0.028019 \\
9-bus  & 317 & 0.684123 & 0.059669 \\
14-bus & 361 & 0.385633 & 0.026457 \\
\hline
\end{tabular}
\end{table}

\begin{figure}[!t]
\centering
\includegraphics[width=\linewidth]{\detokenize{Buses-5_06-07-2026_00-05-09-Obj.png}}
\caption{5-bus objective history from the selected Simbi-only run.}
\label{fig:five-bus-plot}
\end{figure}

\begin{figure}[!t]
\centering
\includegraphics[width=\linewidth]{\detokenize{Buses-9_06-08-2026_19-48-09-Obj.png}}
\caption{9-bus objective history from the selected Simbi-only run.}
\label{fig:nine-bus-plot}
\end{figure}

\begin{figure*}[!t]
\centering
\includegraphics[width=0.95\textwidth]{\detokenize{qhd_vs_standard_summary.png}}
\caption{Summary of closest-to-standard QHD/Simbi iterates, using combined Euclidean distance over bus states and branch flows.}
\label{fig:standard-summary}
\end{figure*}

\subsection{Discussion}
The 14-bus case reaches a maximum residual near $10^{-4}$ and tracks the standard objective within approximately $0.73\%$ at the best-residual iterate. The 5-bus run is close in objective but retains a residual near $10^{-3}$, and the selected 9-bus run is visibly less feasible, with a maximum residual of $7.51\times10^{-3}$. Tables~\ref{tab:load-results} and~\ref{tab:closest-standard} show that the residual-minimizing iterate and the closest-to-standard iterate need not coincide. This is expected in a penalty-based nonconvex method: a point can be close to a reference OPF solution in bus/branch variables while still having larger lifted residuals, or vice versa. These results indicate that the linearized ALM mechanism can produce meaningful ACOPF iterates through a discretized spin backend, but the feasibility quality is sensitive to penalty scaling, stepsize adaptation, unary resolution, and local refinement.

The present study has three important limitations. First, all numerical results in this draft use the Simbi backend rather than quantum hardware. Second, the table does not yet include a systematic comparison against Gurobi, IPOPT, MATPOWER, or other standard ACOPF baselines \needref{MATPOWER/IPOPT/Gurobi ACOPF baseline references}. Third, the encoded spin dimension grows with both the number of lifted variables and unary resolution; future work should study adaptive resolution, warm starts, and better scaling of the ALM penalty.

\section{Conclusion}
This paper presented a lifted rectangular ACOPF formulation and a linearized augmented Lagrangian algorithm designed to generate quadratic box-constrained subproblems. These subproblems are compatible with QHD-style Ising optimizers and were tested with a Simulated Bifurcation backend. Preliminary comparisons on 5-, 9-, and 14-bus systems show that the framework can approach the standard-answer objectives while reducing AC residuals, with the strongest feasibility performance observed on the selected 14-bus run. The immediate next steps are to add systematic classical solver baselines, improve feasibility refinement, report runtime and sensitivity studies, and replace the red reference placeholders with formal citations.

\end{document}
